% Draft: reordered opening of the Introduction (K's arc, 2026-11-24, 13:3x).
% Paragraph 1: tES acts in the brain's frequency range but cannot focus deep (theorem).
% Paragraph 2: TI works with a nonlinear feature of the field that escapes the restriction.
% Paragraph 3: how TI builds it (beat), Eq. (1), envelope definition, need for nonlinearity.
% Replaces the current first three paragraphs (up to "The mechanism has three stages").
% All statements and citations are already in the paper; only the order and links are new.
% "a few hundred hertz": K suggested DC to ~250 Hz; adjust if he prefers a number.

\section{Introduction}
Transcranial electrical stimulation (tES), or transcranial current stimulation (tCS),
injects weak currents through scalp electrodes \cite{ruffini2013} and acts directly at
the frequencies of brain function, from DC to a few hundred hertz
\cite{reato2010,liu2018,krause2023}. Its reach to deep targets is limited
by physics. Under the quasistatic approximation, electric-field strength has no strict
three-dimensional interior maximum in homogeneous, isotropic, source-free tissue: its
maximum lies on a boundary, such as the scalp or a conductivity interface
\cite{heller1992}. No montage can therefore make a deep region of uniform tissue receive
a stronger field than its surroundings. With spatially fixed linear coupling to the
membrane, the same restriction applies to the mean-square polarization and to any other
map that inherits subharmonicity from the field (\S\ref{sec:dc}).

Temporal-interference (TI) stimulation aims at the focal modulation of deep brain
targets by working with a nonlinear feature of the field, its \textit{envelope}
\cite{grossman2017}. The envelope excursion and the quadratic beat amplitude can have
interior maxima (supplement, \S\ref{supp:interior-maxima}), whereas the beat-period mean
of the envelope remains subject to the restriction (\S\ref{sec:dc}). Spatial targeting
therefore depends on which feature of the waveform the neurons detect.

To construct such a feature, two sinusoidal currents are injected at $f_1$ and
$f_2=f_1+\Delta f$ through separate electrode pairs, typically with $f_1\sim1$--$2$~kHz
and $\Delta f\sim1$--$100$~Hz, producing, where they overlap, an amplitude-modulated (AM)
electric field whose direction generally rotates over the beat period. Changing the
currents steers the region of largest envelope modulation. The difference frequency
itself is absent from the input,
\begin{equation}
s(t)=\varepsilon_1\cos(\omega_1 t)+\varepsilon_2\cos(\omega_2 t),\qquad
\omega_{1,2}=2\pi f_{1,2},\quad \Omega=\omega_2-\omega_1=2\pi\,\Delta f .
\label{eq:am}
\end{equation}
Here $s(t)$ is the local membrane polarization, $t$ is time, and
$\varepsilon_{1,2}$ are its tone amplitudes in mV, produced by coupling of the applied
field to the membrane; $\Omega$ is the beat angular frequency, and
$f_c=(f_1+f_2)/2$ is the carrier frequency. The signal \textit{envelope} is the modulus
of the analytic signal, $\mathcal A(t)=|s(t)+i\mathcal H[s](t)|$, where $\mathcal H$ is
the Hilbert transform \cite{gabor1946}. The spectrum of $s(t)$ contains only the two
applied lines. A linear, time-invariant system can attenuate or phase-shift them but
cannot create a new one. Passive membrane filtering therefore cannot demodulate TI
\cite{mirzakhalili2020}; a nonlinear element must mix the carriers into the
low-frequency band, as in AM reception \cite{nahin2001}. Neural targeting thus depends
on how neurons and neural populations transduce this waveform.

% Bibliography entry to add to references.bib (verified against Crossref, 2026-11-24):
% @article{ruffini2013,
%   author  = {Ruffini, Giulio and Wendling, Fabrice and Merlet, Isabelle and Molaee-Ardekani, Behnam and Mekonnen, Abeye and Salvador, Ricardo and Soria-Frisch, Aureli and Grau, Carles and Dunne, Stephen and Miranda, Pedro C.},
%   title   = {Transcranial Current Brain Stimulation ({tCS}): Models and Technologies},
%   journal = {IEEE Transactions on Neural Systems and Rehabilitation Engineering},
%   volume  = {21}, number = {3}, pages = {333--345}, year = {2013},
%   doi     = {10.1109/TNSRE.2012.2200046}}
